\chapter{Step Inputs: asking the Operator for a value}\label{ch:inputs}

A step declares its inputs with \texttt{:setInputs(\{\dots\})}, keyed from 1 in the order they
are asked for, and reads them with \texttt{ATE.input(n)}. Every declaration needs \texttt{desc}
and \texttt{type}; the rest depend on the type.

\begin{center}
\footnotesize
\begin{tabular}{@{}>{\raggedright\arraybackslash}p{1.8cm}>{\raggedright\arraybackslash}p{3.0cm}>{\raggedright\arraybackslash}p{9.4cm}@{}}
\toprule
Field & Applies to & Meaning \\
\midrule
\texttt{desc} & all & the text shown to the Operator (required) \\
\texttt{type} & all & \texttt{ENUM}, \texttt{INTEGER}, \texttt{FLOAT}, \texttt{STRING},
  \texttt{PROMPT} or \texttt{POLL} (required) \\
\texttt{input} & \texttt{ENUM} & the widget: \texttt{RADIOBUT}, \texttt{CHECKBOX} or
  \texttt{PUSHBUT} \\
\texttt{enum} & \texttt{ENUM} & the option labels, at least one \\
\texttt{range} & \texttt{INTEGER}, \texttt{FLOAT} & \texttt{\{min=, max=\}}; empty means
  unbounded \\
\texttt{pattern} & \texttt{STRING} & a regular expression the value must match \\
\texttt{prompt} & \texttt{PROMPT}, \texttt{POLL} & the button text: a string, or an array of
  strings \\
\texttt{tooltip} & all & hover help \\
\texttt{orient} & all & \texttt{VERT} or \texttt{HORIZ} \\
\bottomrule
\end{tabular}
\end{center}

A field used on a type it does not belong to is a load error, so an \texttt{enum} on an
\texttt{INTEGER} is caught when the file loads rather than at run time.

\section{What \texttt{ATE.input(n)} returns}

\begin{itemize}
  \item \texttt{RADIOBUT} or \texttt{PUSHBUT} --- the label the Operator chose, as a string.
  \item \texttt{CHECKBOX} --- a Lua table of every label ticked; empty when none are.
  \item \texttt{INTEGER} --- a Lua integer. \texttt{FLOAT} --- a Lua float.
  \item \texttt{STRING} --- the text as typed.
  \item \texttt{PROMPT} --- blocks the step until the Operator presses the button, then returns
    its text.
  \item \texttt{POLL} --- never blocks. It returns \texttt{nil} until a button is clicked, then
    that button's text, and keeps returning it until the Attempt starts again. A POLL click is
    recorded when it happens, not when the schedule reads it.
\end{itemize}

Asking for an input the step did not declare is an error, and a value outside a range or a
pattern is refused before the step sees it.

\noindent\textbf{\texttt{inputs.tslua}}\par
\lstinputlisting[language=ATELua]{../examples/test-writer/inputs.tslua}

By the end of this chapter you can put a prompt on a step and read the Operator's answer back
in the right type. \texttt{POLL} is the one to reach for when a step has to watch for a click
while it works, and the loop in step 8 is the pattern.

\cando{declare every kind of Step Input a schedule needs, and read each one back without
guessing its type.}
